% ===== PRÉAMBULE CANONIQUE — à coller VERBATIM en tête de CHAQUE .tex =====
\documentclass[11pt,a4paper]{article}
\usepackage[utf8]{inputenc}\usepackage[T1]{fontenc}\usepackage{lmodern}
\usepackage[french]{babel}
\usepackage{amsmath,amssymb,mathtools}\usepackage{icomma}
\usepackage{siunitx}\sisetup{locale=FR}  % \qty{}{}, \si{}, \ang{}
\usepackage{xcolor}\usepackage{colortbl}
\usepackage{tikz}\usepackage{pgfplots}\pgfplotsset{compat=1.18}
\usepackage[siunitx,europeanresistors]{circuitikz} % circuits (élec.)
\usepackage[most]{tcolorbox}\usepackage{enumitem}\usepackage{array,booktabs}
\usepackage[a4paper,margin=2.2cm]{geometry}
\usetikzlibrary{arrows.meta,calc,angles,quotes,positioning,patterns,%
  backgrounds,shapes.geometric,decorations.pathmorphing,decorations.markings,intersections}
\definecolor{primary}{RGB}{222,49,99}\definecolor{primarydark}{RGB}{150,22,55}
\definecolor{defcol}{RGB}{0,114,178}\definecolor{methcol}{RGB}{52,92,148}
\definecolor{excol}{RGB}{0,135,90}\definecolor{attncol}{RGB}{200,40,55}
\tcbset{boxbase/.style={enhanced,breakable,boxrule=0.9pt,arc=2.5pt,
  left=3mm,right=3mm,top=2.3mm,bottom=2.3mm,fonttitle=\bfseries,coltitle=white}}
% --- boîtes TUTORIEL (noms EXACTS, ne pas renommer) ---
\newtcolorbox{cle}[1][]{boxbase,colback=defcol!6,colframe=defcol,
  title={\textsc{À retenir}\ifx\relax#1\relax\relax\else\textmd{ : \itshape #1}\fi}}
\newtcolorbox{methode}[1][]{boxbase,colback=methcol!7,colframe=methcol,sharp corners,
  title={\textsc{Méthode}\ifx\relax#1\relax\relax\else\textmd{ --- \itshape #1}\fi}}
\newtcolorbox{exemplebox}[1][]{boxbase,colback=excol!5,colframe=excol,
  title={\textsc{Exemple}\ifx\relax#1\relax\relax\else\textmd{ : \itshape #1}\fi}}
\newtcolorbox{piege}[1][]{boxbase,colback=attncol!6,colframe=attncol,
  title={\textsc{Piège}\ifx\relax#1\relax\relax\else\textmd{ : \itshape #1}\fi}}
\newtcolorbox{remarque}[1][]{boxbase,colback=black!4,colframe=black!45,
  title={\textsc{Remarque}\ifx\relax#1\relax\relax\else\textmd{ : \itshape #1}\fi}}
% --- boîtes EXERCICE / CORRIGÉ (noms EXACTS, auto-comptées) ---
\newtcolorbox[auto counter]{exo}[1][]{boxbase,colback=primary!4,colframe=primary,
  title={\textsc{Exercice}~\thetcbcounter\ifx\relax#1\relax\relax\else\textmd{ --- \itshape #1}\fi}}
\newtcolorbox[auto counter]{corrige}[1][]{boxbase,colback=excol!4,colframe=excol,
  title={\textsc{Corrigé}~\thetcbcounter\ifx\relax#1\relax\relax\else\textmd{ --- \itshape #1}\fi}}
\newcommand{\vv}[1]{\vec{#1}}\newcommand{\ds}{\displaystyle}
\newcommand{\R}{\mathbb{R}}\newcommand{\N}{\mathbb{N}}\newcommand{\Z}{\mathbb{Z}}
% =========================================================================
\begin{document}
% THEME_TITLE: Indice, célérité et longueur d'onde

\begin{center}
{\LARGE\bfseries\color{primarydark} Thème 3 — Indice, célérité, longueur d'onde}\\[2pt]
{\color{primary}\itshape Ce qui change (et ce qui ne change pas) d'un milieu à l'autre}
\end{center}

\begin{cle}[l'indice de réfraction]
L'indice $n_i$ d'un milieu mesure à quel point il « ralentit » la lumière :
\[ \boxed{\,n_i=\dfrac{c}{c_i}\,}\qquad\Longleftrightarrow\qquad c_i=\dfrac{c}{n_i} \]
avec $c\approx\qty{3.0e8}{m/s}$ la célérité dans le vide et $c_i$ la célérité dans le milieu. Comme
$c_i\le c$, on a toujours $n_i\ge1$. \textbf{Plus $n$ est grand, plus la lumière est lente.}
\end{cle}

\begin{cle}[la relation d'onde et l'invariant fréquence]
Pour une onde : $\boxed{\,c_i=\lambda_i\,f\,}$, soit $\lambda_i=c_i/f$.

Lorsque la lumière change de milieu, la \textbf{fréquence $f$ reste constante} (elle est imposée par
la source). Ce sont la \emph{célérité} $c_i$ et la \emph{longueur d'onde} $\lambda_i$ qui changent.
\begin{center}\itshape « $f$ fixe ; $c$ et $\lambda$ variables. »\end{center}
\end{cle}

\begin{cle}[la chaîne des rapports]
En combinant $n_i=c/c_i$, $\lambda_i=c_i/f$ ($f$ constante) et la loi de Snell–Descartes, tous les
rapports s'enchaînent :
\[ \boxed{\;\dfrac{\sin\hat{\imath}_1}{\sin\hat{\imath}_2}=\dfrac{n_2}{n_1}=\dfrac{c_1}{c_2}
   =\dfrac{\lambda_1}{\lambda_2}\;} \]
\end{cle}

\begin{methode}[relier deux grandeurs entre deux milieux]
\begin{enumerate}[leftmargin=1.7em,topsep=2pt,itemsep=1pt]
  \item Écrire la chaîne et \textbf{repérer les deux rapports} utiles (angles, indices, célérités
        ou longueurs d'onde).
  \item Attention à l'\textbf{inversion} : les indices $n_2/n_1$ sont dans le sens inverse des
        célérités $c_1/c_2$ (grand indice $\Rightarrow$ petite vitesse).
  \item Les longueurs d'onde suivent les célérités : $\lambda_1/\lambda_2=c_1/c_2$.
  \item Pour une longueur d'onde dans un milieu depuis celle dans le vide : $\lambda_i=\lambda_0/n_i$.
\end{enumerate}
\end{methode}

\begin{exemplebox}[air $\to$ verre, ce qui change]
Une lumière verte ($f=\qty{5.0e14}{Hz}$) passe de l'air ($n_1=1$) au verre ($n_2=1{,}5$) :
\[ c_1=\frac{c}{1}=\qty{3.0e8}{m/s},\quad c_2=\frac{c}{1{,}5}=\qty{2.0e8}{m/s}, \]
\[ \lambda_1=\frac{c_1}{f}=\qty{600}{nm},\quad \lambda_2=\frac{c_2}{f}=\qty{400}{nm},\quad f\ \text{inchangée}. \]
La vitesse et la longueur d'onde ont été divisées par $1{,}5$ ; la fréquence (donc la couleur) est
restée la même.
\end{exemplebox}

\begin{piege}[ne pas inverser les rapports]
Un \emph{grand} indice correspond à une \emph{petite} vitesse : c'est pourquoi $n_2/n_1$ et $c_1/c_2$
sont « à l'envers » l'un de l'autre. En revanche $\lambda$ suit $c$ (même sens). Vérifiez toujours par
le bon sens : dans un milieu plus réfringent, la lumière est plus \emph{lente} et sa longueur d'onde
plus \emph{courte}.
\end{piege}
\end{document}
